
\subsubsection{3.1 Overview}

The methodology comprises three components:

\begin{enumerate}
\item \textbf{Inductive Bias Measurement:} Quantify locality versus global attention through coefficient of variation (CoV) of layer-wise weight updates during training
\item \textbf{Architecture-Task Interaction:} Test whether DWS excels on local anomaly detection while NFT excels on holistic property aggregation
\item \textbf{Controlled Comparison:} Match training procedures across architectures to isolate inductive bias effects
\end{enumerate}

\subsubsection{3.2 Architectures}

\textbf{MLP Baseline:} A multi-layer perceptron that flattens weight matrices into vectors without equivariant structure. Configuration: 3 hidden layers with ReLU activations, approximately 9.8M parameters.

\textbf{Deep Weight Space (DWS):} Processes weights through equivariant layers that respect permutation symmetry while preserving spatial structure within each layer. Configuration: 3 equivariant layers, hidden dimension 256, approximately 4.9M parameters.

\textbf{Neural Functional Transformer (NFT):} Tokenizes weight matrices and applies transformer self-attention across all tokens. Configuration: 4 attention layers, 4 heads, hidden dimension 256, approximately 5.3M parameters.

Note: Parameter counts are not matched within the originally specified 10\% tolerance (DWS 4.9M, NFT 5.3M, MLP 9.8M). This introduces a potential capacity confound, though the MLP with approximately twice the parameters of the equivariant methods does not achieve superior performance on the accuracy prediction task.

\subsubsection{3.3 Inductive Bias Measurement}

The coefficient of variation (CoV) of layer-wise weight updates quantifies locality versus global processing:

$$\text{CoV} = \frac{\sigma(\|\Delta W_l\|)}{\mu(\|\Delta W_l\|)}$$

where $\Delta W_l$ is the weight update magnitude for layer $l$.

Higher CoV indicates more varied updates across layers (locality preservation); lower CoV indicates more uniform updates (global information sharing).

\subsubsection{3.4 Tasks}

\textbf{Backdoor Detection (Local Pattern):} Binary classification of synthetic weight populations as clean or containing localized perturbations. Metric: Area Under ROC Curve (AUC).

\textbf{Accuracy Prediction (Global Statistic):} Regression to predict a synthetic target computed from global weight statistics (sigmoid of mean norms and means across layers). Metric: Root Mean Squared Error (RMSE).

Both tasks use synthetic weight populations derived from MNIST Implicit Neural Representation (INR) models, with 600 training samples and 200 test samples.

\subsubsection{3.5 Experimental Design}

A $2\times 3$ factorial design (Task $\times$ Architecture) with two-way ANOVA tests for interaction effects. Training configuration: AdamW optimizer, learning rate 1e-4, weight decay 1e-2, batch size 64, 100 epochs for h-m3 experiments, seeds [42, 123, 456].
